\section{Questão 15}

Considerando a condições de contorno apresentadas no experimento randômico descrito na questão 10:

% \begin{multicols}{2}
\begin{itemize}
    \item [a)] 
    calcule a expectância condicional $\E{X\mid Y = y}$;
    
    \item [b)] 
    calcule a expectância condicional $\E{Y\mid X = x}$;
    
    \item [c)]
    Expresse as expectâncias condicionais calculadas em a) e b), $\E{X\mid Y}$ e $\E{Y\mid X}$, respectivamente, como variáveis randômicas cujo valores são representados por 
    $g(y) = \E{X\mid Y = y}$ quando $Y = y$ e 
    $g(x) = \E{Y\mid X = x}$ quando $X = x$;
    
    \item [d)]
    calcule a variância condicional $\Var{X\mid Y=y}$;
    
    \item [e)]
    calcule a variância condicional $\Var{Y\mid X=x}$;
    
    \item [f)]
    expresse as variâncias condicionais $\Var{X\mid Y=y}$ e $\Var{Y\mid X=x}$ como variáveis randômicas.
\end{itemize}
% \end{multicols}

\begin{respostas}
    % O procedimento para obter as respostas são análogos aos já apresentados nas questões 11 a 14.
    
    \begin{itemize}
        \item [a), b)] 
        % calcule a expectância condicional $\E{X\mid Y = y}$;
        % calcule a expectância condicional $\E{Y\mid X = x}$;
        \begin{minipage}{0.48\linewidth}
            \centering
            \begin{tabular}{c cccc}
            \toprule
                $y$ & 0 & 1 & 2 & 3 \\
            \midrule
                $\E{X\mid Y}$ & $6/4$ & $6/4$ & $6/4$ & $6/4$ \\
            \bottomrule
            \end{tabular}
        \end{minipage}
        %
        \begin{minipage}{0.48\linewidth}
            \centering
            \begin{tabular}{c cccc}
            \toprule
                $x$ & 0 & 1 & 2 & 3 \\
            \midrule
                $\E{Y\mid X}$ & $6/4$ & $6/4$ & $6/4$ & $6/4$ \\
            \bottomrule
            \end{tabular}
        \end{minipage}
        
        \item [c)]
        $g(z) = 3/2$ se $z\in\{0, 1, 2, 3\}$ e $g(z)=0$ caso contrário.
        
        \item [d), e)]
        % calcule a variância condicional $\Var{X\mid Y=y}$.
        % calcule a variância condicional $\Var{Y\mid X=x}$.
        % calcule a expectância condicional $\E{X\mid Y = y}$;
        % calcule a expectância condicional $\E{Y\mid X = x}$;
        \begin{minipage}{0.48\linewidth}
            \centering
            \begin{tabular}{c cccc}
            \toprule
                $y$ & 0 & 1 & 2 & 3 \\
            \midrule
                $\Var{X\mid Y}$ & $5/4$ & $5/4$ & $5/4$ & $5/4$ \\
            \bottomrule
            \end{tabular}
        \end{minipage}
        %
        \begin{minipage}{0.48\linewidth}
            \centering
            \begin{tabular}{c cccc}
            \toprule
                $x$ & 0 & 1 & 2 & 3 \\
            \midrule
                $\Var{Y\mid X}$ & $5/4$ & $5/4$ & $5/4$ & $5/4$ \\
            \bottomrule
            \end{tabular}
        \end{minipage}
        
        \item [f)]
        % Expresse as variâncias condicionais $\Var{X\mid Y=y}$ e $\Var{Y\mid X=x}$ como variáveis randômicas.
        $h(z) = 5/4$ se $z\in\{0, 1, 2, 3\}$ e $h(z)=0$ caso contrário.
    \end{itemize}

\end{respostas}